\subsection{Biểu đồ 1: Phân phối giá niêm yết theo borough}

\subsubsection{Thiết kế Idiom}

\begin{table}[H]
\centering
\renewcommand{\arraystretch}{1.1}
\begin{tabularx}{\textwidth}{|l|l|X|}
\hline
\textbf{Idiom} & \multicolumn{2}{l|}{\textbf{Box-and-Whisker Plot}} \\ \hline
\textbf{What}
& \multicolumn{2}{X|}{Borough (neighbourhood\_group\_cleansed): Categorical} \\ \cline{2-3}
& \multicolumn{2}{X|}{Price: Quantitative} \\ \cline{2-3}
& \multicolumn{2}{X|}{Room Type: Categorical} \\ \hline
\textbf{How}
& \textbf{Encode} &
\textbf{Mark:} Glyph (Area cho thân hộp, Line cho râu/trung vị, Point cho ngoại lệ) \newline
\textbf{Channel:} \newline
\hspace*{1em} Pos X: Borough (Categorical -- phân biệt 5 borough) \newline
\hspace*{1em} Pos Y: Price (Quantitative -- IQR, median, whiskers) \newline
\hspace*{1em} Hue Color: Room Type (Categorical -- 4 loại phòng) \\ \cline{2-3}
& \textbf{Manipulate} &
Selection: Hover để xem chi tiết (median, Q1, Q3, min, max, count). \\ \cline{2-3}
& \textbf{Reduce} &
Filter: price\_is\_outlier = False (loại extreme outliers bằng IQR method). \\ \hline
\textbf{Why} & \multicolumn{2}{X|}{
\textbf{produce $\rightarrow$ compare $\rightarrow$ discover} \newline
So sánh phân phối giá đầy đủ giữa 5 borough. Glyph là idiom duy nhất hiển thị Q1, median, Q3, whiskers đồng thời --- bar chart chỉ cho median, histogram chỉ cho 1 nhóm.
} \\ \hline
\textbf{Scale} & \multicolumn{2}{X|}{
Main key: 5 (borough). \newline
Color key: 4 (room type). \newline
Items: $\sim$21,000 listing (sau lọc).
} \\ \hline
\end{tabularx}
\end{table}

\begin{figure}[H]
    \centering
    \safeincludegraphics[width=0.9\linewidth]{charts/domain_task5_chart1.png}
    \caption{Hình 5.1. Glyph phân phối giá niêm yết theo borough tại NYC}
    \label{fig:domain-task5-chart1}
\end{figure}

\subsubsection{Đánh giá biểu đồ}

\textbf{1. Nguyên lý biểu đạt (Expressiveness):}

\begin{itemize}
    \item Thuộc tính: Quận / Borough -- neighbourhood\_group\_cleansed (C), Giá niêm yết / Price (Q), Loại phòng / Room Type (C), price\_is\_outlier (C -- filter).
    \item Channel:
    \begin{itemize}
        \item PosX: phân biệt 5 borough (C) $\rightarrow$ phù hợp. Spatial region là channel hiệu quả nhất cho categorical attribute.
        \item PosY: thể hiện phạm vi phân phối giá (Q1, Median, Q3, whiskers) $\rightarrow$ phù hợp Quantitative. Position on common scale là channel chính xác nhất cho quantitative.
        \item Hue Color: phân biệt loại phòng (C) $\rightarrow$ phù hợp. Số màu = 4 $\leq$ 7 $\rightarrow$ nằm trong giới hạn discriminability của HUE channel.
    \end{itemize}
    \item Đánh giá mức độ quan trọng:
    \begin{itemize}
        \item Ưu tiên 1 -- Vị trí dọc (PosY --- Position on common scale): PosY là kênh quan trọng nhất trong Glyph. Kênh này mã hóa toàn bộ phân phối giá (Quantitative) của từng borough: giá trị Q1, Median, Q3, whiskers và outlier đều được biểu diễn theo thang đo chung trên trục dọc. Theo Mackinlay, Position on a common scale là kênh hiệu quả nhất cho dữ liệu định lượng, cho phép người xem ước lượng và so sánh khoảng cách giá giữa các borough với độ chính xác cao nhất --- đây là nền tảng để thực hiện task compare (so sánh phân phối) và discover (phát hiện outlier).
        \item Ưu tiên 2 -- Vị trí ngang (PosX --- Position / Spatial region): PosX phân tách 5 borough thành các nhóm rời nhau theo chiều ngang, mã hóa thuộc tính Borough (Categorical/Nominal) theo không gian. Đây là kênh quan trọng thứ hai vì nó thiết lập cấu trúc so sánh chính của biểu đồ: người xem định vị từng borough theo vị trí ngang trước khi đọc độ cao cột. Thứ tự cột được sắp xếp giảm dần theo giá median giúp việc so sánh diễn ra tức thì và trực quan hơn.
        \item Ưu tiên 3 -- Màu sắc (Color Hue): Color Hue mã hóa Room Type (Categorical/Nominal) bằng màu sắc phân biệt cho từng loại phòng. Theo Mackinlay, Color Hue phù hợp nhất cho dữ liệu danh mục vì mỗi màu được cảm nhận là ngang hàng nhau, không ngụ ý thứ tự hay mức độ. Kênh này cho phép người xem đồng thời so sánh phân phối giá giữa các loại phòng trong cùng một borough mà không cần tách thành nhiều biểu đồ riêng biệt --- tăng mật độ thông tin mà không làm tăng độ phức tạp hiển thị, với điều kiện số màu $\leq$ 7 như hiện tại.
    \end{itemize}
    \item Kết luận: Các channel đúng với bản chất dữ liệu. Glyph là idiom chuẩn mực nhất để biểu diễn phân phối liên tục theo nhóm --- đặc biệt khi cần so sánh IQR chứ không chỉ median.
\end{itemize}

\textbf{2. Nguyên lý hiệu quả (Effectiveness):}

\begin{itemize}
    \item Độ chính xác (Accuracy): PosY (Position on common scale) là channel chính xác nhất theo channel effectiveness ranking --- xếp hạng 1 trong hierarchy. Hue Color dùng cho categorical không liên quan đến sai số định lượng.
    \item Khả năng phân biệt (Discriminability): 4 màu cho 4 room type trong phạm vi $\leq$ 7 $\rightarrow$ mắt người phân biệt hoàn toàn tốt. Whiskers và box IQR hiển thị rõ ràng trên nền trắng.
    \item Khả năng tách biệt (Separability): PosX, PosY và Color tách biệt hoàn toàn. Filter price\_is\_outlier = False loại bỏ extreme outliers giúp chart dễ đọc hơn.
\end{itemize}

\subsubsection{Phân tích biểu đồ (Insight)}

\begin{itemize}
    \item Manhattan có median price cao nhất ($\sim$\$150--200) và IQR rộng nhất --- tồn tại cả phân khúc bình dân lẫn cao cấp trong cùng một borough.
    \item Bronx và Staten Island có IQR hẹp và median thấp ($\sim$\$92--99) --- thị trường giá ổn định, phù hợp du khách cần ngân sách dự đoán được.
    \item Loại phòng Entire home/apt luôn có box cao hơn Private room ở mọi borough --- sự phân tầng giá theo loại phòng nhất quán trên toàn thành phố.
    \item Khuyến nghị: Du khách ngân sách trung bình nên cân nhắc Brooklyn --- vị trí gần Manhattan nhưng median giá thấp hơn đáng kể.
\end{itemize}
